\subsection{Recognition without supplied call boundaries}\label{sec:no-call}
\textbf{Withdrawing call endpoints tests a different attacker representation, while whole-episode association remains available.}
We reanalyze all 3,000 recorded episodes using complete encrypted connections, including initialization and closing records. The attacker receives only relative record times, directions, ciphertext lengths and destination identities. It receives no call starts/ends, operation names, role labels or HTTP stream assignment. A per-connection idle-gap rule proposes events; five gaps from 5 to 250\,ms compete on validation task accuracy. Raw-record classifiers bypass this eventization entirely.

For each condition, LR and RF use destination summaries and early/late event sequences, with and without timing. A training-only GMM removes short-gap events before a classifier trained on no-cover traces. Validation separately chooses the baseline and removal attacks. This is a retrospective, fixed-protocol analysis of reused fixtures; only the immediate-arm paired comparison is a primary test, with Holm over the two models. Table~\ref{tab:no-call-main} also gives each finite suite's largest test count without substituting that maximum into the primary test.

\begin{table}[t]\centering\small
\caption{No supplied call boundaries. All encrypted records in associated episode connections enter the analysis. Cells give selected idle-grid baseline / removal / largest test count including the later exploratory protocol-aware suite, out of 120; uniform reference20. Whole-task association remains supplied.}\label{tab:no-call-main}
\begin{tabularx}{\columnwidth}{@{}Yrr@{}}\toprule
Policy & Qwen & Llama \\\midrule
No cover & 115/109/117 & 108/101/112 \\
Immediate & 80/69/91 & 90/86/97 \\
History, sequential & 73/74/75 & 88/56/93 \\
History, concurrent & 56/61/62 & 78/62/88 \\
Planner shadow & 69/65/73 & 69/56/69 \\
\bottomrule\end{tabularx}\end{table}

On immediate cover, selected removal recognizes 69/120 Qwen and 86/120 Llama tasks, against baselines 80 and 90. The paired differences are -9.2 and -3.3 points (Holm $p=0.02545$ and $0.5413$). Removal underperforms the baseline in Qwen after correction; Llama does not reject. The original removal advantage is not reproduced by this boundary-free idle-gap pipeline. After observing the idle-grid results, an exploratory protocol-aware eventizer uses the public request-padding quantum and MCP setup sequence: large upward records mark proposed requests, and large downward records delimit responses before the next request. It does not read call truth. Within this representation, selected removal changes Qwen recognition from 68 to 91, but Llama from 90 to 87; these exploratory counts do not replace the original tests. Immediate cover permits maxima 91/120 and 97/120 across the record-only suites. Whole-episode association remains given, and differing representations prevent interpreting the change from the earlier call-window count as a pure causal boundary effect. Supplement S3 gives the full protocol and independent event-matching diagnostic.

The independent event audit gives protocol-event F1 0.295 and 0.380 at IoU .5 on immediate cover. Task recognition from record summaries therefore does not imply accurate recovery of calls.
